% Nouveau EGA V (cote 25): archive cote 25, site page 23 (PDF page 24).
% typescript. Draft compared: none.
% AI-transcribed by Claude Opus 5.5 (claude-opus-5-5, adaptive thinking, effort high), Message Batches API, batch ega25b (msgbatch_01A7n5YuXubUPNK1zsepi8RG), 2026-11-13; not yet checked by eye.
\scanpage{25/23}
\mspage{\hand{20}}

\uline{Proposition} 5.3.   Avec les notations précédentes, soit U l'ensemble des
$x\in X_{\hand{\check{}}P}$ , d'image $\hand{\zeta}$ dans $\hand{\check{}}P$ , tels que $\phi_{\hand{\zeta}}$ soit $F_{k(\hand{\zeta})}$-régulière en x. \deleted{Suppo-%
sons de plus} \inserted{Si} F de présentation finie et plat relativement à S\hand{,} Alors:
% NOTE: the words "Si F de présentation finie et plat relativement à S," are encircled by hand, and a line leads from the loop down to a), between "a)" and "U".
a) U est ouvert, et $\underline{G}\hand{|}U$ est plat relativement à $\hand{\check{}}P$ .

b) Pour tout $s\in\hand{\check{}}P$ , si $\hand{\eta}$ désigne le point générique de $\hand{\check{}}P_s$ , U con-%
tient $X_{k(\hand{\eta})}$.

Démonstration. a) Comme $F_{\hand{\check{}}P}$ est \inserted[typed]{de présentation finie, et} plat par rapport à $\hand{\check{}}P$ , la conclu-%
sion résulte de 11.3. (et même de $\underline{E}_{III}$ ... dans le cas S loc noeth).

b) On peut supposer S=Spec(k). Les cycles associés à $F_{k(\hand{\eta})}$ sont
\inserted[typed]{\xout{\illegible{1}}}
(en vertu du par.3 ) les images inverses des cycles \inserted{$Z_i$} associés à F. \xout{En
vertu de 2.6, par exemple, si $Z_i$ est une partie \unsure{sous-schéma} fermé intègre
de X telle que} Si $f(Z_i)$ est fini, alors en vertu de 2.3. \xout{X}
$Z_{i\,k(\hand{\eta})}\cap Y = \emptyset$ , dans le cas contraire, on vertu de 2.6.\inserted[typed]{par exemple} on a
également $Z_{i\,k(\hand{\eta})}\cap Y \neq Z_{i\,k(\hand{\eta})}$ (raisons de dimensions; 2.3. déjà invo-%
qué dans 2.6. est s ans doute une meulleure raison), ce qui prouve que $\phi_{\hand{\eta}}$
ne s'annule sur aucun des $Z_{i\,k(\hand{\eta})}$ , et prouve donc b).

\uline{Corollaire} 5.4. Soit V l'ensemble des $\hand{\zeta}\in\hand{\check{}}P$ tels que $\phi_{\hand{\zeta}}$ soit \xout{E} $F_{k(\hand{\zeta})}$-régu-%
lier. \inserted{Si F est de présentation finie} \over{A}{\unsure{a}}lors V est constructible et contient les points génériques des
fibres de $\hand{\check{}}P$ sur S, d'autre part si \inserted{de plus} X est propre sur S\inserted{et F plat sur S}\hand{,} l'ensemble V est
ouvert. \deleted{\hand{Si \unsure{$S$ spec $(A)$} (\unsure{$A$ complet}) alors \unsure{en} tous cas $\unsure{f}(\check{P}-V)$ est
de codimension $\unsure{\geq|\neq} 2$.}}

\inserted[typed]{\uline{Remarque}}\xout{Xxxxxxxxxx}5.5. \xout{Supposons} soit $\hand{\zeta}\in\hand{\check{}}P$ au dessus de $s\in S$, et supposons
que $F_s$ soit sans cycles associés immergés, alors on voit aussitôt que
$\hand{\zeta}\in V$ (notation de 5.4.) signifie aussi que\xout{xxxxx} aucune composante
irréductible de supp $F_{k(\hand{\zeta})}$ ne se trouve au dessus de $H_{\hand{\zeta}}$ (et on peut
évidemment dans ce critère remplacer $k(\hand{\zeta})$ par une extension arbitraire
de $k(\hand{\zeta})$ ). Notons que l'hypothèse $(S_1)$ sur $F_s$ qu'on vient de faire es

\deleted{\hand{$\times$ Ce dernier fait se vérifie par la méthode de 5.3. b) et la \unsure{théorie}
de la dimension (on remplace $\eta$ par $\zeta$ \unsure{tel que} $\mathcal{O}_{\check{P},\zeta}$ de dim $1$)}}
\margin{23}
